\chapter{Hướng dẫn cài đặt và chạy hệ thống}
\label{app:install}

\section{Yêu cầu môi trường}

\begin{itemize}
  \item Node.js v18+, npm v9+
  \item Python 3.9+
  \item PostgreSQL 12+ (production) hoặc SQLite (development)
  \item Docker và Docker Compose (để chạy container)
\end{itemize}

\section{Cài đặt Backend}

\begin{lstlisting}[language=bash, caption={Cài đặt và chạy Backend}]
cd WEB/multilanguage_transolator_be

# Tao moi truong ao
python -m venv env
source env/bin/activate    # Linux/Mac
# env\Scripts\activate     # Windows

# Cai dat thu vien
pip install -r requirements.txt

# Cau hinh bien moi truong
cp .env.example .env
# Chinh sua .env: SECRET_KEY, DB_*, AWS_*, GOOGLE_APPLICATION_CREDENTIALS, ...

# Chay migration
python manage.py migrate

# Khoi dong server
python manage.py runserver   # http://127.0.0.1:8000
\end{lstlisting}

\section{Cài đặt Frontend}

\begin{lstlisting}[language=bash, caption={Cài đặt và chạy Frontend}]
cd WEB/multilanguage_transolator_fe

# Cai dat thu vien
npm install

# Khoi dong development server
npm run dev    # http://localhost:5173
\end{lstlisting}

\section{Biến môi trường quan trọng}

\begin{table}[H]
\centering
\caption{Biến môi trường Backend}
\renewcommand{\arraystretch}{1.3}
\begin{tabularx}{\textwidth}{|l|X|}
\hline
\textbf{Biến}                         & \textbf{Mô tả} \\
\hline
\texttt{SECRET\_KEY}                   & Django secret key \\
\hline
\texttt{DEBUG}                         & True (dev) / False (prod) \\
\hline
\texttt{DB\_NAME, DB\_USER, DB\_PASSWORD, DB\_HOST, DB\_PORT} & PostgreSQL connection \\
\hline
\texttt{AWS\_ACCESS\_KEY\_ID}           & AWS credentials \\
\hline
\texttt{AWS\_SECRET\_ACCESS\_KEY}       & AWS credentials \\
\hline
\texttt{AWS\_STORAGE\_BUCKET\_NAME}     & S3 bucket name \\
\hline
\texttt{GOOGLE\_APPLICATION\_CREDENTIALS} & Đường dẫn file JSON service account Google \\
\hline
\texttt{PROJECT\_ID}                   & Google Cloud project ID \\
\hline
\texttt{INPUT\_URI}                    & GCS URI cho glossary CSV (gs://bucket/file.csv) \\
\hline
\texttt{JWT\_ACCESS\_TTL\_MINUTES}      & TTL access token (mặc định: 15) \\
\hline
\texttt{FRONTEND\_URL}                 & URL frontend (CORS) \\
\hline
\end{tabularx}
\end{table}

\chapter{Đặc tả Use Case chi tiết}
\label{app:usecase}

\section{Đặc tả UC-05: Xem lịch sử dịch}

\begin{table}[H]
\centering
\caption{Đặc tả UC-05: Xem lịch sử dịch}
\renewcommand{\arraystretch}{1.3}
\begin{tabularx}{\textwidth}{|l|X|}
\hline
\textbf{Tên use case}    & UC-05: Xem lịch sử dịch \\
\hline
\textbf{Tác nhân}        & Người dùng đã đăng nhập \\
\hline
\textbf{Điều kiện trước} & Người dùng đã thực hiện ít nhất một lần dịch \\
\hline
\textbf{Mô tả}           & Người dùng xem danh sách tất cả tệp đã dịch, có thể lọc, tìm kiếm
                           và tải lại tệp kết quả \\
\hline
\textbf{Luồng chính}     & Người dùng truy cập \texttt{/file-history} → Hệ thống trả về danh
                           sách \texttt{TranslatedFile} của user hiện tại, sắp xếp theo
                           \texttt{created\_at} giảm dần → Người dùng có thể click URL để
                           tải tệp về từ S3 \\
\hline
\end{tabularx}
\end{table}
